\documentclass[10pt,a4paper,twocolumn]{article}
\usepackage[T1]{fontenc}
\usepackage{lmodern}
\usepackage[margin=2cm,columnsep=0.7cm]{geometry}
\usepackage{amsmath,amssymb}
\usepackage{graphicx}
\usepackage{booktabs}
\usepackage[backend=biber,style=numeric,sorting=none]{biblatex}
\addbibresource{refs.bib}
\usepackage[hidelinks]{hyperref}
\usepackage{cleveref}

\title{Title of the Paper}
\author{First Author\thanks{Example University} \and Second Author\thanks{Example Institute}}
\date{\today}

\begin{document}

% The optional argument of \twocolumn sets the title block and abstract across both columns.
\twocolumn[
  \begin{@twocolumnfalse}
    \maketitle
    \begin{abstract}
      Summarise the paper in about 150 words: problem, method, main result.
    \end{abstract}
    \vspace{1em}
  \end{@twocolumnfalse}
]

\section{Introduction}
State the problem and what you contribute. Cite with \verb|\cite|, for example~\cite{knuth84}.
Narrow columns need short paragraphs and compact figures.

\section{Method}\label{sec:method}
Define the approach. A numbered equation:
\begin{equation}
  \hat{\theta} = \arg\min_\theta \sum_{i=1}^{n} \ell\bigl(f_\theta(x_i), y_i\bigr).
  \label{eq:erm}
\end{equation}
\Cref{eq:erm} is referred to with \verb|\cref|.

\section{Results}
Use \texttt{figure*} and \texttt{table*} for floats that need the full page width.

\begin{table}[t]
  \centering
  \caption{A one-column table.}
  \begin{tabular}{lrr}
    \toprule
    Method & Accuracy & Time (s) \\
    \midrule
    Baseline & 0.81 & 12 \\
    Ours & 0.87 & 15 \\
    \bottomrule
  \end{tabular}
\end{table}

\section{Conclusion}
Wrap up in a few sentences.

\printbibliography

\end{document}
